%==============================================================================
%  RESUME BASE — TRACK B  (education-forward layout)
%
%  Section order: Summary -> Education (with relevant modules) -> Projects
%                 -> Experience -> Skills -> Certifications -> Achievements
%  Use for: graduates, career changers, and any track where coursework and
%  projects carry more weight than employment history.
%
%  Same preamble and commands as Track A, so a resume can be switched between
%  bases without reformatting. Same rules apply: never alter the preamble,
%  never leave a {{TOKEN}} in a sent resume, escape % & _ # $ in content.
%==============================================================================


\documentclass[letterpaper,10.5pt]{article}

\usepackage{latexsym}
\usepackage[empty]{fullpage}
\usepackage{titlesec}
\usepackage{marvosym}
\usepackage[usenames,dvipsnames]{color}
\usepackage{verbatim}
\usepackage{enumitem}
\usepackage{xcolor}
\usepackage[
    colorlinks=true,
    linkcolor=black,   % internal links stay black
    urlcolor=blue,     % external links in blue
    citecolor=blue
]{hyperref}   % loaded exactly once - a second \usepackage{hyperref} with different
              % options is an 'Option clash' error
\usepackage{fancyhdr}
\usepackage[english]{babel}
% fontawesome5 is optional: guarded so the file still builds on a minimal TeX install
\IfFileExists{fontawesome5.sty}{\usepackage{fontawesome5}}{}
\usepackage{tabularx}
%   - Keep the glyphtounicode line and \pdfgentounicode=1 — they make the PDF's
%     text extractable by ATS parsers.
\usepackage{ifthen}
\pagestyle{fancy}
\fancyhf{}
\fancyfoot{}
\renewcommand{\headrulewidth}{0pt}
\renewcommand{\footrulewidth}{0pt}

% Adjust margins tighter
\addtolength{\oddsidemargin}{-0.55in}
\addtolength{\evensidemargin}{-0.55in}
\addtolength{\textwidth}{1.1in}
\addtolength{\topmargin}{-0.8in}
\addtolength{\textheight}{1.35in}

\urlstyle{same}

\raggedbottom
\raggedright
\setlength{\tabcolsep}{0in}
\setlength{\parskip}{0pt} % remove paragraph gaps

% Compact lists
\setlist[itemize]{topsep=1pt, itemsep=0.5pt, parsep=0pt, leftmargin=0.15in}

% Section formatting
\titleformat{\section}{
  \vspace{2pt}\scshape\raggedright\normalsize
}{}{0em}{}[\color{black}\titlerule \vspace{1pt}]

\pdfgentounicode=1

%-------------------------
% Custom commands
\newcommand{\resumeItem}[1]{
  \item\small{{#1}}
}

\newcommand{\resumeSubheading}[4]{
  \vspace{1pt}\item
    \begin{tabular*}{0.97\textwidth}[t]{l@{\extracolsep{\fill}}r}
      \textbf{#1} & #2 \\
      \textit{\small#3} & \textit{\small #4} \\
    \end{tabular*}\vspace{0.5pt}
}

\newcommand{\resumeProjectHeading}[2]{
    \item
    \begin{tabular*}{0.97\textwidth}{l@{\extracolsep{\fill}}r}
      \small#1 & #2 \\
    \end{tabular*}\vspace{1pt}
}

\newcommand{\resumeSubItem}[1]{\resumeItem{#1}\vspace{-1pt}}

\renewcommand\labelitemii{$\vcenter{\hbox{\tiny$\bullet$}}$}

\newcommand{\resumeSubHeadingListStart}{\begin{itemize}[leftmargin=0.15in, label={}]} 
\newcommand{\resumeSubHeadingListEnd}{\end{itemize}}
\newcommand{\resumeItemListStart}{\begin{itemize}[leftmargin=0.15in]}
\newcommand{\resumeItemListEnd}{\end{itemize}\vspace{0.5pt}}

% -------------------------
% Hierarchical Experience Commands
% -------------------------

% Robust empty-argument test (expansion-safe, no ifthen quirks).
\newcommand{\ifnotempty}[2]{\if\relax\detokenize{#1}\relax\else#2\fi}

% \job{Employer}{subtitle or empty}{extra or empty}{date, right-aligned}
% Empty arguments are skipped cleanly - no "There's no line here to end" error.
\newcommand{\job}[4]{%
    \vspace{2pt}%
    \noindent\textbf{#1}\hfill #4\par
    \ifnotempty{#2}{{\small #2}\par}%
    \ifnotempty{#3}{#3\par}%
    \vspace{2pt}%
}

% \project{Role or project title}{location or empty}{extra or empty}
\newcommand{\project}[3]{%
    \vspace{1pt}%
    \noindent\textit{#1}\ifnotempty{#2}{ #2}\par
    \ifnotempty{#3}{#3\par}%
    \vspace{1pt}%
}

\newcommand{\subproject}[1]{%
    \vspace{0.5pt}%
    \noindent\textbf{#1}\par
    \vspace{-1pt}%
}

% Passthrough wrapper used around the Experience section content.
% (Defined here so the document compiles cleanly with no "undefined control sequence" error.)
\newcommand{\sectioncontent}[1]{#1}

% soul gives the underlined links; guarded with a plain-underline fallback so the
% file still builds on a minimal TeX install (Overleaf always has soul)
\IfFileExists{soul.sty}{\usepackage{soul}}{}
\providecommand{\ul}[1]{\underline{#1}}  % fallback if soul is unavailable

%-------------------------------------------
%%%%%%  RESUME STARTS HERE  %%%%%%%%%%%%%%%%%%%%%%%%%%%%
\begin{document}

%----------HEADING----------
\begin{center}
    \textbf{\Huge \scshape Annie Prasanna Manoharan} \\ \vspace{1pt}
    \small {{PHONE}} $|$
    \href{mailto:{{EMAIL}}}{\textcolor{blue}{\ul{{{EMAIL}}}}} $|$
    \href{{{LINKEDIN_URL}}}{\textcolor{blue}{\ul{LinkedIn}}} $|$
    \href{{{GITHUB_URL}}}{\textcolor{blue}{\ul{GitHub}}} \\
    {{CITY}}, United States $|$ {{WORK_AUTH}}
\end{center}

%-----------SUMMARY-----------
% Graduate formula: <degree + specialisation> + <strongest project or coursework
% area matching the JD> + <the domain interest this role sits in>.
% Never claim years of experience that do not exist.
\section{Summary}
\begin{itemize}[leftmargin=0.15in,label={}]
    \item \small{ {{SUMMARY_PARAGRAPH}} }
\end{itemize}

%-----------EDUCATION-----------
% Education leads on this base. Include the modules the JD actually names.
\section{Education}
  \resumeSubHeadingListStart
    \resumeSubheading
      {{{INSTITUTION_1}}}{{{EDU_LOCATION_1}}}
      {{{DEGREE_1}}}{{{EDU_DATES_1}}}
      \resumeItemListStart
        \resumeItem{\textbf{Relevant modules:} {{EDU_COURSEWORK_1}}}
        \resumeItem{\textbf{Major project:} {{EDU_PROJECT_1}}}
      \resumeItemListEnd
    \resumeSubheading
      {{{INSTITUTION_2}}}{{{EDU_LOCATION_2}}}
      {{{DEGREE_2}}}{{{EDU_DATES_2}}}
      \resumeItemListStart
        \resumeItem{\textbf{Relevant modules:} {{EDU_COURSEWORK_2}}}
      \resumeItemListEnd
  \resumeSubHeadingListEnd

%-----------PROJECTS-----------
% On this base the projects do the work employment normally does: three of them,
% each with what it does, the stack, and the candidate's own contribution.
\section{Projects}
\job{{{PROJECT_B1_NAME}}}{}{}{\hspace{1mm}{{PROJECT_B1_DATES}}}

\begin{itemize}
    \item {{PROJECT_B1_DESCRIPTION}}
    \item \textbf{Stack:} {{PROJECT_B1_STACK}}
    \item {{PROJECT_B1_CONTRIBUTION}}
\end{itemize}

\job{{{PROJECT_B2_NAME}}}{}{}{\hspace{1mm}{{PROJECT_B2_DATES}}}

\begin{itemize}
    \item {{PROJECT_B2_DESCRIPTION}}
    \item \textbf{Stack:} {{PROJECT_B2_STACK}}
    \item {{PROJECT_B2_CONTRIBUTION}}
\end{itemize}

\job{{{PROJECT_B3_NAME}}}{}{}{\hspace{1mm}{{PROJECT_B3_DATES}}}

\begin{itemize}
    \item {{PROJECT_B3_DESCRIPTION}}
    \item {{PROJECT_B3_CONTRIBUTION}}
\end{itemize}

%-----------EXPERIENCE----------- (internships, part-time, volunteer — labelled honestly)
\section{Experience}
\sectioncontent{

\job{{{EMPLOYER_1}}}{}{}{\hspace{1mm}{{EMP_DATES_1}}}

    \project{{{ROLE_1}} ({{EMP_TYPE_1}})}{{{EMP_LOCATION_1}}}{}
    \begin{itemize}
        \item {{EMP_BULLET_1_1}}
        \item {{EMP_BULLET_1_2}}
    \end{itemize}

}
% If there is no employment history at all, delete this entire section rather
% than relabelling coursework or personal projects as jobs.

%-----------TECHNICAL SKILLS-----------
% Keep coursework-level knowledge in its own line, described as familiarity.
\section{Technical Skills}
\begin{itemize}[leftmargin=0.15in, label={}]
    \small{\item{
        \textbf{ {{SKILL_CATEGORY_1}} }{: {{SKILL_LIST_1}} } \\
        \textbf{ {{SKILL_CATEGORY_2}} }{: {{SKILL_LIST_2}} } \\
        \textbf{ {{SKILL_CATEGORY_3}} }{: {{SKILL_LIST_3}} } \\
        \textbf{ {{SKILL_CATEGORY_4}} }{: {{SKILL_LIST_4}} } \\
        \textbf{Academic exposure}{: {{SKILL_LIST_EXPOSURE}} }
    }}
\end{itemize}

%-----------CERTIFICATIONS-----------
\section{Certifications}
\begin{itemize}[leftmargin=0.15in]
    \item \small{ {{CERT_1}} }
    \item \small{ {{CERT_2}} }
\end{itemize}

%-----------ACHIEVEMENTS----------- (delete if empty)
\section{Achievements}
\begin{itemize}[leftmargin=0.15in]
    \item \small{ {{AWARD_1}} }
\end{itemize}

\end{document}
